\documentclass{article}
\usepackage{amsmath,amssymb,amsthm}
\usepackage{enumitem}
\usepackage{hyperref}
\usepackage{geometry}
\geometry{margin=1in}
\newtheorem{theorem}{Theorem}
\newtheorem{lemma}{Lemma}
\newtheorem{assumption}{Assumption}
\newcommand{\EE}{\mathbb{E}}
\newcommand{\RR}{\mathbb{R}}
\newcommand{\norm}[1]{\left\|#1\right\|}
\newcommand{\inner}[1]{\left\langle #1 \right\rangle}
\begin{document}

\section*{Algorithm Name}
\textbf{Differentially Private Stochastic Gradient Descent (DP‑SGD) with Gaussian Noise}

\section*{Mathematical Formulation}
Let $f:\RR^{d}\rightarrow\RR$ be a convex, $L$‑smooth objective function.  
At iteration $t=0,1,\dots ,T-1$ the algorithm performs

\[
\begin{aligned}
\text{(1)}\qquad &\tilde{g}_t \;=\; \nabla f(x_t) \;+\; \xi_t, \\[2mm]
\text{(2)}\qquad &x_{t+1}\;=\; x_t \;-\; \eta_t \tilde{g}_t ,
\end{aligned}
\]

where  

\begin{itemize}
    \item $\eta_t>0$ is the learning‑rate (possibly constant $\eta_t=\eta$);
    \item $\xi_t\sim\mathcal N(0,\sigma^2 I_d)$ is an independent Gaussian noise vector,  
          introduced to satisfy a given $(\varepsilon,\delta)$‑differential privacy budget.
\end{itemize}

For convenience we denote the (noisy) gradient by $g_t:=\tilde{g}_t$.

\section*{Pseudocode}
\begin{verbatim}
Input: initial point x0, learning rate schedule {ηt}, noise scale σ,
       total iterations T
for t = 0,…,T‑1 do
    g = ∇f(x_t)                     // true gradient
    ξ ~ N(0, σ^2 I_d)               // Gaussian noise
    g̃ = g + ξ                      // noisy gradient
    x_{t+1} = x_t - η_t * g̃        // SGD update
end for
return x_T
\end{verbatim}

\section*{Dry Run Example}
Consider the one‑dimensional quadratic $f(w)=(w-3)^2$, thus $\nabla f(w)=2(w-3)$.  
Choose $w_0=0$, constant step size $\eta=0.1$, and set $\sigma=0$ (no privacy noise) for illustration.

\begin{center}
\begin{tabular}{c|c|c|c}
$t$ & $w_t$ & $\nabla f(w_t)$ & $w_{t+1}=w_t-\eta\nabla f(w_t)$ \\ \hline
0 & $0$ & $2(0-3)=-6$ & $0-0.1(-6)=0.6$ \\
1 & $0.6$ & $2(0.6-3)=-4.8$ & $0.6-0.1(-4.8)=1.08$ \\
2 & $1.08$ & $2(1.08-3)=-3.84$ & $1.08-0.1(-3.84)=1.464$ \\
\end{tabular}
\end{center}

Corresponding objective values: $f(0)=9$, $f(0.6)=5.76$, $f(1.08)=3.6864$, $f(1.464)=2.3660$,
showing descent toward the optimum $w^\star=3$.

\section*{Assumptions}
\begin{assumption}[Convexity]\label{as:convex}
$f$ is convex: $\forall x,y\in\RR^d,\; f(y)\ge f(x)+\inner{\nabla f(x)}{y-x}$.
\end{assumption}
\begin{assumption}[L‑smoothness]\label{as:smooth}
$f$ has $L$‑Lipschitz gradients: $\forall x,y,\;
\norm{\nabla f(x)-\nabla f(y)}\le L\norm{x-y}$.
\end{assumption}
\begin{assumption}[Bounded Gradient Norm]\label{as:bound}
There exists $G>0$ such that $\norm{\nabla f(x)}\le G$ for all $x$ in the domain visited by the algorithm.
\end{assumption}
\begin{assumption}[Gaussian Noise]\label{as:noise}
$\{\xi_t\}_{t\ge0}$ are i.i.d. $\mathcal N(0,\sigma^2 I_d)$, independent of $x_t$.
Hence $\EE[\xi_t]=0$ and $\EE[\|\xi_t\|^2]=\sigma^2 d$.
\end{assumption}
\begin{assumption}[Step‑size]\label{as:stepsize}
The learning rates are non‑increasing and satisfy $\eta_t\le \frac{1}{L}$ for all $t$.
\end{assumption}

\section*{Main Convergence Theorem}
\begin{theorem}\label{thm:main}
Assume \ref{as:convex}–\ref{as:stepsize} hold. Let $x^\star$ be a minimizer of $f$ and define the averaged iterate
\[
\bar x_T \;=\;\frac{1}{T}\sum_{t=0}^{T-1} x_t .
\]
For a constant stepsize $\eta_t=\eta\le\frac{1}{L}$, the expected suboptimality satisfies
\[
\EE\!\bigl[f(\bar x_T)-f(x^\star)\bigr]
\;\le\;
\frac{\norm{x_0-x^\star}^2}{2\eta T}
\;+\;
\frac{\eta}{2}\bigl(G^2+\sigma^2 d\bigr).
\tag{3}
\]
Choosing $\eta=\frac{\norm{x_0-x^\star}}{\sqrt{T\,(G^2+\sigma^2 d)}}$ (which respects $\eta\le 1/L$ for sufficiently large $T$) yields
\[
\EE\!\bigl[f(\bar x_T)-f(x^\star)\bigr]
\;\le\;
\frac{\norm{x_0-x^\star}\sqrt{G^2+\sigma^2 d}}{\sqrt{T}} .
\tag{4}
\]
Thus DP‑SGD attains the standard $O(1/\sqrt{T})$ rate, with an additional variance term caused by the privacy noise.
\end{theorem}

\section*{Theoretical Properties}
\begin{enumerate}[label=\arabic*.]
\item \textbf{Stability.} The recursion \eqref{eq:recursion} (derived below) shows that the distance $\|x_{t+1}-x^\star\|$ never explodes as long as $\eta_t\le1/L$; the Gaussian term only adds a bounded variance $\sigma^2 d$.
\item \textbf{Hyper‑parameter Sensitivity.}
    \begin{itemize}
        \item Larger $\eta$ speeds up the deterministic part (the $G^2$ term) but amplifies the noise contribution $\eta\sigma^2 d$.
        \item The optimal $\eta$ balances the two, leading to the choice in \eqref{4}.
    \end{itemize}
\item \textbf{Comparison to non‑private SGD.} Setting $\sigma=0$ recovers the classical bound
    $\EE[f(\bar x_T)-f(x^\star)]\le \frac{\norm{x_0-x^\star}^2}{2\eta T}+\frac{\eta G^2}{2}$,
    i.e., the privacy‑induced variance term is the only extra price.
\end{enumerate}

\section*{Discussion}
The theorem demonstrates that adding isotropic Gaussian noise for differential privacy does not destroy the $O(1/\sqrt{T})$ convergence guarantee of SGD on convex smooth problems. The privacy budget influences $\sigma$; a tighter $(\varepsilon,\delta)$ requirement forces larger $\sigma$, which appears linearly under the square‑root in the final bound. Consequently, to keep the same accuracy one must increase the number of iterations $T$ (or equivalently the dataset size) proportionally to $\sigma^2 d$.

\section*{Preliminaries (Definitions and Lemmas)}
\begin{lemma}[Three‑point identity for smooth functions]\label{lem:smooth}
If $f$ is $L$‑smooth, then for any $x,y$,
\[
f(y) \le f(x) + \inner{\nabla f(x)}{y-x} + \frac{L}{2}\norm{y-x}^2 .
\tag{5}
\]
\end{lemma}
\begin{proof}
Standard consequence of the mean‑value theorem for vector‑valued functions; see e.g. \cite{Nesterov2004}. \hfill $\square$
\end{proof}
\begin{lemma}[Unbiasedness of noisy gradient]\label{lem:unbiased}
For the noisy gradient $g_t=\nabla f(x_t)+\xi_t$, we have
\[
\EE[g_t\mid x_t]=\nabla f(x_t),\qquad
\EE\bigl[\|g_t\|^2\mid x_t\bigr]=\|\nabla f(x_t)\|^2+\sigma^2 d .
\tag{6}
\]
\end{lemma}
\begin{proof}
Since $\EE[\xi_t]=0$ and $\xi_t$ is independent of $x_t$, the first identity follows immediately.  
For the second,
\[
\EE\bigl[\|g_t\|^2\mid x_t\bigr]
= \EE\bigl[\|\nabla f(x_t)+\xi_t\|^2\mid x_t\bigr]
= \|\nabla f(x_t)\|^2 + 2\inner{\nabla f(x_t)}{\EE[\xi_t]} + \EE\|\xi_t\|^2
= \|\nabla f(x_t)\|^2 + \sigma^2 d .
\hfill $\square$
\]
\end{proof}

\section*{Main Proof (Step‑by‑Step Derivation)}
We prove Theorem~\ref{thm:main}.  
All expectations are taken with respect to the randomness of the noise sequence $\{\xi_t\}$.

\bigskip
\noindent\textbf{Step 1. Distance recursion.}  
From the update rule $x_{t+1}=x_t-\eta g_t$,
\[
\begin{aligned}
\|x_{t+1}-x^\star\|^2
&= \|x_t-\eta g_t - x^\star\|^2  \\
&= \|x_t-x^\star\|^2 -2\eta\inner{g_t}{x_t-x^\star}
   +\eta^2\|g_t\|^2 .
\end{aligned}
\tag{7}
\]
\noindent\textbf{[Step 1]}

\bigskip
\noindent\textbf{Step 2. Take conditional expectation.}  
Condition on $x_t$ and use Lemma~\ref{lem:unbiased}:
\[
\EE\!\left[\|x_{t+1}-x^\star\|^2\mid x_t\right]
= \|x_t-x^\star\|^2
 -2\eta\inner{\nabla f(x_t)}{x_t-x^\star}
 +\eta^2\bigl(\|\nabla f(x_t)\|^2+\sigma^2 d\bigr).
\tag{8}
\]
\noindent\textbf{[Step 2]}

\bigskip
\noindent\textbf{Step 3. Use convexity.}  
By Assumption~\ref{as:convex},
\[
f(x_t)-f(x^\star) \le \inner{\nabla f(x_t)}{x_t-x^\star}.
\tag{9}
\]
Insert (9) into (8) and rearrange:
\[
\begin{aligned}
2\eta\bigl(f(x_t)-f(x^\star)\bigr)
&\le \|x_t-x^\star\|^2 - \EE\!\left[\|x_{t+1}-x^\star\|^2\mid x_t\right] \\
&\qquad +\eta^2\bigl(\|\nabla f(x_t)\|^2+\sigma^2 d\bigr) .
\end{aligned}
\tag{10}
\]
\noindent\textbf{[Step 3]}

\bigskip
\noindent\textbf{Step 4. Bound the gradient norm.}  
Using the bounded‑gradient assumption \ref{as:bound},
$\|\nabla f(x_t)\|^2\le G^2$, thus
\[
2\eta\bigl(f(x_t)-f(x^\star)\bigr)
\le \|x_t-x^\star\|^2 - \EE\!\left[\|x_{t+1}-x^\star\|^2\mid x_t\right]
   +\eta^2\bigl(G^2+\sigma^2 d\bigr).
\tag{11}
\]
\noindent\textbf{[Step 4]}

\bigskip
\noindent\textbf{Step 5. Take full expectation and telescope.}  
Let $\Delta_t:=\EE\|x_t-x^\star\|^2$. Taking expectation of (11) yields
\[
2\eta\,\EE\!\bigl[f(x_t)-f(x^\star)\bigr]
\le \Delta_t - \Delta_{t+1}
   +\eta^2\bigl(G^2+\sigma^2 d\bigr).
\tag{12}
\]
Sum (12) over $t=0,\dots,T-1$:
\[
2\eta\sum_{t=0}^{T-1}\EE\!\bigl[f(x_t)-f(x^\star)\bigr]
\le \Delta_0 - \Delta_T + T\eta^2\bigl(G^2+\sigma^2 d\bigr).
\tag{13}
\]
Since $\Delta_T\ge0$, we drop it and divide by $2\eta T$:
\[
\frac{1}{T}\sum_{t=0}^{T-1}\EE\!\bigl[f(x_t)-f(x^\star)\bigr]
\le
\frac{\Delta_0}{2\eta T}
     +\frac{\eta}{2}\bigl(G^2+\sigma^2 d\bigr).
\tag{14}
\]
\noindent\textbf{[Step 5]}

\bigskip
\noindent\textbf{Step 6. From iterate average to function average.}  
By convexity,
$f(\bar x_T)\le\frac1T\sum_{t=0}^{T-1}f(x_t)$. Taking expectation and applying (14) gives exactly bound \eqref{3} of Theorem~\ref{thm:main}.

\noindent\textbf{[Step 6]}

\bigskip
\noindent\textbf{Step 7. Optimizing the stepsize.}  
The right‑hand side of \eqref{3} is
\[
\Phi(\eta)=\frac{\|x_0-x^\star\|^2}{2\eta T}+\frac{\eta}{2}(G^2+\sigma^2 d).
\]
Minimizing w.r.t.\ $\eta>0$ yields $\eta^\star
= \frac{\|x_0-x^\star\|}{\sqrt{T\,(G^2+\sigma^2 d)}}$, and substituting $\eta^\star$ into $\Phi(\eta)$ gives bound \eqref{4}.

\noindent\textbf{[Step 7]}

Thus the theorem is proved. \hfill $\blacksquare$

\section*{Rate Derivation (With Constants)}
From \eqref{4} we obtain the explicit constant:
\[
\EE\!\bigl[f(\bar x_T)-f(x^\star)\bigr]
\le
\underbrace{\norm{x_0-x^\star}}_{\text{initial distance}}
\;\underbrace{\sqrt{G^2+\sigma^2 d}}_{\text{combined variance}}
\;T^{-1/2}.
\]
If $\sigma=0$, the term reduces to $\norm{x_0-x^\star}G\,T^{-1/2}$, the classical SGD rate.

\section*{Special Cases}
\begin{enumerate}[label=\alph*.]
\item \textbf{No privacy noise ($\sigma=0$).} The bound collapses to the standard convex‑smooth SGD guarantee.
\item \textbf{Strongly convex $f$.} If $f$ is $\mu$‑strongly convex, the same analysis (with Lemma~\ref{lem:smooth} replaced by the strong‑convexity inequality) yields a linear convergence term plus a steady‑state error of order $\eta\sigma^2 d/\mu$.
\item \textbf{Decreasing stepsize $\eta_t=\frac{c}{\sqrt{t+1}}$.} A similar telescoping argument gives
\[
\EE\!\bigl[f(\bar x_T)-f(x^\star)\bigr]\le
O\!\Bigl(\frac{1}{\sqrt{T}}\Bigr)
\]
with constants involving $\sum_{t=0}^{T-1}\eta_t$ and $\sum_{t=0}^{T-1}\eta_t^2$.
\end{enumerate}

\begin{thebibliography}{9}
\bibitem{Nesterov2004}
Y.~Nesterov,
\textit{Introductory Lectures on Convex Optimization: A Basic Course},
Springer, 2004.
\end{thebibliography}

\end{document}